\documentclass[10pt,a4paper]{amsart}
\usepackage[margin=19mm]{geometry}
\usepackage{amsmath,amssymb,mathtools}
\usepackage[scr=rsfs,cal=euler]{mathalfa}
\usepackage{xfrac}
\usepackage[unicode,hidelinks,bookmarks=false]{hyperref}
\newcommand{\ie}{i.e.\ }
\newcommand{\cf}{cf.\ }
\newcommand{\resp}{respectively\ }
\newcommand{\Subsection}{\subsection}
\providecommand{\nobreakdash}{\nobreak\hbox{-}}
\setlength{\emergencystretch}{2em}
\multlinegap0pt
\usepackage[shortlabels]{enumitem}
\usepackage{tikz}
\usepackage{xcolor}
\usepackage{xfrac}
\usepackage{amssymb}
\usepackage[normalem]{ulem}
\newtheorem{proposition}{Proposition}
\newtheorem{hypothesis}{Hypothesis}
\newtheorem{definition}{Definition}
\newtheorem{lemma}{Lemma}
\newcommand{\R}{\mathbb R}
\newcommand{\al}{\alpha}
\newcommand{\beq}{\begin{equation}}
\newcommand{\bzeta}{{\zeta}}
\newcommand{\cac}{\mathcal C}
\newcommand{\ck}{\mathcal K}
\newcommand{\cn}{\mathcal N}
\newcommand{\cp}{\mathcal P}
\newcommand{\cq}{\mathcal Q}
\newcommand{\cv}{\mathcal V}
\newcommand{\der}{\delta}
\newcommand{\eeq}{\end{equation}}
\newcommand{\ga}{\gamma}
\newcommand{\iot}{\int_{0}^{t}}
\newcommand{\ka}{\kappa}
\newcommand{\lp}{\left(}
\newcommand{\rp}{\right)}
\newcommand{\si}{\sigma}
\title{The Pontryagin maximum principle and $Q$-functions in rough environments}
\author{Estepan Ashkarian, Prakash Chakraborty, Harsha Honnappa, Samy Tindel}
\date{Source: arXiv:2601.05354v1 (CC BY 4.0)}
\begin{document}
\maketitle

\noindent\emph{Source and licence.} The Pontryagin maximum principle and $Q$-functions in rough environments. Version arXiv:2601.05354v1 (CC BY 4.0), distributed by arXiv under the Creative Commons Attribution 4.0 International licence, http://creativecommons.org/licenses/by/4.0/, as recorded in the arXiv licence metadata for this exact version and shown on the abstract page https://arxiv.org/abs/2601.05354v1. Full licence notice: \texttt{licenses/arxiv-2601.05354-CC-BY-4.0.txt}.\par\smallskip

\noindent\textit{Editorial scope.} Counted unit: the article's lemma X (source0.tex lines 850--867). The article's complete original proof of it is delivered with it (source0.tex lines 869--1085, one \texttt{proof} environment of 217 lines). No mathematics is written, repaired or re-derived: the statement and the proof are the article's own lines, in the article's own notation, and no retained line is substituted or re-flowed.  Retained uncounted as the source-local context the proof needs: lines 248--250; lines 336--343; lines 371--384; lines 396--408; lines 475--492; lines 585--600; lines 589--591; lines 602--615; lines 611--613; lines 660--662; lines 807--813; lines 827--844; lines 841--842.  Nothing was omitted from any retained span, so the package carries no omission record.  The retained lines cite 3 works, whose entries are transcribed verbatim from the article's own bibliography source in the e-print and delivered with short editorial labels recorded in the item's reference mapping.  No retained line contains a figure environment, an image-inclusion command or drawing code, so no image is delivered and the package declares no asset dependency.  Editorial formatting only, in addition: an AMS \texttt{amsart} preamble that restates the article's own declarations and macros used by the retained lines, namely the environments \texttt{proposition}, \texttt{hypothesis}, \texttt{definition}, \texttt{lemma} on separate counters, together with the standard packages those lines need.  The retained environments print in this document as Proposition 1, Hypothesis 1, Definition 1, Lemma 1 in order of appearance, whereas the article prints the same items with the numbering of its own full document.  The complete native source of the article as distributed in the arXiv e-print for this exact version is bundled unchanged as \texttt{sources/192333/source0.tex}.
\begin{equation}\label{d1}
x_t^{\ga} = x_0 + \iot \int_{U} b(s, x_s^{\ga}, a) \, \ga_s(da) ds + \iot \si(s, x_s^{\ga}) \, d\zeta_s \, ,
\end{equation}

The building block of the {rough path} theory is the so-called sewing map lemma. We recall this fundamental result here.\begin{proposition}\label{prop:La}
	Let $h \in \cac_3^{\mu}(V)$ for $\mu > 1$ be such that $\der h = 0$. Then there exists a unique $g = \Lambda(h) \in \cac_2^{\mu}(V)$ such that $\der g = h$. Furthermore for such an $h$, the following relations hold true:
	\begin{equation*}
	\der \Lambda(h) = h ,
	\quad\text{  and  }\quad
{\|\Lambda h\|}_{\mu} \leq \dfrac{1}{2^{\mu} - 2} {\|h\|}_{\mu} .
	\end{equation*}
\end{proposition}

From now on we will focus on a space $V$ of the form $V = \R^d$. The main assumption on the $\R^d$-valued noise $\bzeta$ of equation~\eqref{d1} or~\eqref{d3} is that it gives rise to a geometric rough path. This assumption can be summarized as follows.
\begin{hypothesis}\label{hyp:zeta}
	The path $\zeta:[0,T] \to \R^d$ belongs to the H\"older space ${\cac}^{\al}([0,T];\R^d)$ with $\al > \frac{1}{3}$ and $\zeta_0 = 0$. In addition $\zeta$ admits a L\'evy area above itself, that is, there exists a two index map $\mathbf{\zeta}^2 : {\mathcal{S}_{2}\left([0,T]\right)} \to \R^{d,d}$ which belongs to $\cac_{2}^{2\al}(\R^{d,d})$ and such that 
	\beq\label{eq:zeta^2}
	\der {\mathbf{\zeta}}_{sut}^{2;ij} = {\der \zeta}_{su}^{i} \otimes {\der \zeta}_{ut}^{j},
	\quad\text{  and  }\quad
	{\mathbf{\zeta}_{st}^{2;ij} + \mathbf{\zeta}_{st}^{2;ji}} =  \der \zeta_{st}^{i} \otimes \der \zeta_{st}^{j} .
	\eeq
	The $\al$-H\"older norm of $\zeta$ is denoted by:
	\beq\label{eq:zeta-norm}
	\|\mathbf{\zeta}\|_{\al} = \cn(\zeta;\cac_1^{\al}([0,T],\R^d))
	+\cn(\mathbf{\zeta}^2;\cac_2^{2\al}([0,T],\R^{d,d})).
	\eeq
\end{hypothesis}

\begin{definition}\label{def:weakly-ctrld}
	Let $z$ be a process in $\cac_{1}^{\ka}(\R^n)$ with $\ka \leq \al$ and $2\ka +\ga > 1$. We say that $z$ is weakly controlled by $\zeta$ if $\der z \in \cac_2^{\ka}(\R^n)$ can be decomposed into
	\begin{equation}\label{eq:weakly-controlled}
	\der z^{i} = z^{\zeta;i i_1} \der \zeta^{i_1} + \rho^{i}, 
	\quad\text{   i.e.  }\quad
{\der z}_{st}^{i} = z_{s}^{\zeta;i i_1}{\der \zeta}_{st}^{i_1} + \rho_{st}^{i} ,
	\end{equation}
	for all $(s,t) \in \mathcal{S}_2 \left([0,T]\right)$ and $i=1, \ldots,n$. In the previous formula we assume $z^{\zeta} \in \cac_{1}^{\ka}(\R^{n,d})$ and $\rho$ is a more regular remainder such that $\rho \in \cac_{2}^{2\ka}(\R^{n})$. The space of weakly controlled paths on $[a,b]$ with H\"older continuity $\ka$ and driven by $\zeta$ will be denoted by $\mathcal{Q}_{\zeta}^{\ka}(\R^n)$, and a process $z \in \mathcal{Q}_{\zeta}^{\ka}(\R^n)$ can be considered as a couple $(z, z^\zeta)$. A possible norm on $\mathcal{Q}_{\zeta}^{\ka}(\R^n)$ is given by
	\begin{multline}\label{eq:norm-controlled}
	\cn[z;\mathcal{Q}_{\zeta}^{\ka}(\R^n)] = 
	\cn[z; \cac_1^{\infty}(\R^n)]+\cn[z; \cac_1^{\ka}(\R^n)] + \cn[z^\zeta; \cac_1^{\infty}(\R^{n,d})] \\+ \cn[z^\zeta; \cac_1^{\ka}(\R^{n,d})] + \cn[\rho; \cac_2^{\ka}(\R^n)].
	\end{multline}
\end{definition}

\begin{definition}\label{def:str-contr}
Similarly to Definition~\ref{def:weakly-ctrld}, consider a process $\nu \in \cac_1^{\ka}(\R^n)$ with $\ka \leq \al$ and $2\ka+\al > 1$. Let us suppose that the increment $\der \nu$ lies in $\cac_2^{\ka}$ and can be decomposed as 
\beq\label{eq:str-contr-a}
\der \nu^j = \nu^{\zeta;j i_1} \der \zeta^{i_1} + \nu^{\bzeta^{2}; j i_1 i_2} \bzeta^{2; i_2 i_1} + \rho^{\nu;j},
\eeq
where the regularity for the paths $\nu^{\zeta}$, $\nu^{\bzeta^2}$ and $\rho^{\nu}$ are respectively
\beq\label{eq:str-contr-b}
\nu^{\zeta}, \nu^{\bzeta^2} \in \cac_1^{\ka}([a,b]), \quad\text{and}\quad \rho^{\nu} \in \cac_2^{3\ka}([a,b]).
\eeq
In addition, we assume that $\nu^{\zeta}$ is also a weakly controlled process as in Definition~\ref{def:weakly-ctrld}, with decomposition
\beq\label{eq:str-contr-l}
\der \nu^{\zeta; ji_1} = \nu^{\zeta^{2}; j i_1 i_2} \der \zeta^{i_2} + \rho^{\nu^{\zeta};j},
\eeq
where $\nu^{\zeta^{2}; j i_1 i_2}$ is like in \eqref{eq:str-contr-a} and $\rho^{\nu^{\zeta}; j}$ is a remainder in $\cac_2^{2\ka}([a,b])$. Then we say that $\nu$ is a strongly controlled process in $[a,b]$. We call $\tilde{\cq}_{\zeta}^{\ka}(\R^n)$ this set of paths, equipped with the semi-norm
$$
\cn[\nu; \tilde{\cq}_{\zeta}^{\ka}(\R^n)] = \cn[\nu;\cac_1^{\ka}] + \cn[\nu^{\zeta};\cac_1^{\infty}] + \cn[\nu^{\zeta^{2}}; \cq_{\zeta}^{\ka}] + \cn[\rho;\cac_2^{3\ka}].
$$
\end{definition}

\subsection{Rough differential equations with relaxed controls}
This section is devoted to recall some of the results concerning controlled equations derived in \cite{CHT}. This will lay the ground for our developments throughout the article. Let us then recall that one starts from a rough path $\bzeta$ as in Hypothesis~\ref{hyp:zeta}. 
%One also considers a compact subset $\cv \subset\subset \R^d$ and a measure-valued (relaxed) control $\ga$ such that
As far as the relaxed control is concerned, our state space will be the following:
\begin{definition}\label{def:rel-control}
Consider a domain $U \subset \R^d$. Then the state space for all the relaxed controls below will be a compact subset $\ck$ of $\cp(U)$, where $\cp(U)$ denotes the set of probability measures on $U$. Notice that one can take $\ck = \cp(U)$ whenever $U$ is a compact domain of $\R^d$.
\end{definition}
We now state our basic assumptions on relaxed controls, which will prevail throughout the article.
\begin{hypothesis}\label{hyp:rel-control-1}
In the sequel we consider some measure-valued paths $\{ \ga_s; 0 \leq s \leq T \}$ such that
\begin{enumerate}[wide, labelwidth=!, labelindent=0pt, label=\textnormal{(\arabic*)}]
\setlength\itemsep{.05in}
\item For every $s$, we have $\gamma_{s}\in \mathcal{K}$, where $\mathcal{K}$ is introduced in Definition \ref{def:rel-control}.
\item The application $s\mapsto \gamma_{s}$ is measurable from $[0,T]$ to $\mathcal{K}$.
\end{enumerate}
\end{hypothesis}

\begin{definition}
Consider a domain $U \subset \R^d$. Then the state space for all the relaxed controls below will be a compact subset $\ck$ of $\cp(U)$, where $\cp(U)$ denotes the set of probability measures on $U$. Notice that one can take $\ck = \cp(U)$ whenever $U$ is a compact domain of $\R^d$.
\end{definition}

\noindent
Next for an initial condition $y \in \R^m$, an initial time $s \in [0,T]$, a relaxed control $\ga$ and our rough path $\bzeta$, let $x^{s,y,\ga}$ be the rough differential equation defined by
\beq\label{eq:relaxed-rde}
x_t^{s,y,\ga} = y+\int_s^t \int_U b(r, x_r^{s,y,\ga},a)\ga_r(da)dr + \int_s^t \si(r, x_r^{s,y,\ga}) \, d \bzeta_r .
\eeq
In \eqref{eq:relaxed-rde}, the assumptions on the coefficients $b$ and $\si$ are rather standard within the rough path theory. They are summarized as follows:
\begin{hypothesis}\label{hyp:coeff}
	We assume that the coefficients $b$ and $\sigma$ in \eqref{eq:relaxed-rde} satisfy $b \in \mathcal C_b^{0,1,1}([0,T] \times \R^m \times U; \R^m)$ and $\sigma \in \mathcal C_b^{1,3}([0,T]\times \R^m; \R^{m,d})$. \emph{Furthermore $b(u, x, a)$ is uniformly continuous in $u$, and it is also assumed to be Lipschitz in space uniformly in time and control.}\end{hypothesis}
We now turn to a slight elaboration of Proposition 3.3 from \cite{CHT}, giving existence and uniqueness for RDEs with relaxed controls.
\begin{proposition}\label{ExistUniq2}
Consider $\al > \frac{1}{3}$ and a rough path $\bzeta$ fulfilling Hypothesis~\ref{hyp:zeta}. The coefficients $b$ and $\si$ are such that Hypothesis~\ref{hyp:coeff} is satisfied. In addition, $\ga \colon [0,T]\to \mathcal{P}(U)$ is a measure-valued process fulfilling Hypothesis \ref{hyp:rel-control-1}. Then equation~\eqref{eq:relaxed-rde} admits a unique solution in the space $\cq_{\bzeta}^{\ka}$ introduced in Definition~\ref{def:str-contr}, for all $\frac{1}{3} < \ka < \ga$.
\end{proposition}
\begin{proof}
    Our proposition establishes existence and uniqueness under weaker assumptions on the control $\gamma$ and the drift term $b$ (compared to Proposition 3.3 and Corollary 3.11 from \cite{CHT}). More specifically, the H\"older assumptions on $s\mapsto \gamma_{s}$ assumed in \cite{CHT} is avoided here. We will need this variant for the perturbative approach developed in the upcoming section. The interested reader can check that the proof of Proposition 3.3 and Corollary 3.11 in \cite{CHT} are unchanged under those weaker hypotheses. We skip the details for the sake of brevity.\end{proof}

\begin{proposition}
Consider $\al > \frac{1}{3}$ and a rough path $\bzeta$ fulfilling Hypothesis~\ref{hyp:zeta}. The coefficients $b$ and $\si$ are such that Hypothesis~\ref{hyp:coeff} is satisfied. In addition, $\ga \colon [0,T]\to \mathcal{P}(U)$ is a measure-valued process fulfilling Hypothesis \ref{hyp:rel-control-1}. Then equation~\eqref{eq:relaxed-rde} admits a unique solution in the space $\cq_{\bzeta}^{\ka}$ introduced in Definition~\ref{def:str-contr}, for all $\frac{1}{3} < \ka < \ga$.
\end{proposition}

\begin{equation}\label{d3}
\partial_t v(t,y) + \sup_{\gamma \in \mathcal P(U)} \left\{ \nabla v(t,y) \cdot \int b(t,y,a)\gamma(da) + F(t,y,\gamma) \right\} + \nabla v(t,y) \cdot \sigma(t,y)  \dot{\zeta}_t = 0 \, ,
\end{equation}

        \begin{definition} \label{def:spike variation}
          Suppose that $(\overline{x},\overline{\gamma})$ is an optimal pair. Consider a time variable $t_{0}\in[0,T)$. For an additional variable $\beta>0$, define the interval $I=I_{\beta}=[t_{0},t_{0}+\beta]$ and assume $I\subset [0,T]$. Then we define the spike variation $\gamma^{\beta}$ of the optimal relaxed control as
          \begin{equation}\label{spike}\gamma^{\beta}_{t}=\begin{cases} \mu \hspace{0.2cm} \text{ if } \hspace{0.2cm} t\in I\\
			    \overline{\gamma}_{t} \hspace{0.2cm} \text{ if } \hspace{0.2cm} t\notin I
			\end{cases}\end{equation}
            where $\mu$ is a constant probability measure sitting in $\ck$ (recall that $\ck$ is the compact set of $\mathcal{P}(U)$ introduced in Definition~\ref{def:rel-control}).
           \end{definition}

\begin{proposition}\label{prop: Flow and Jacobian}
    Let us suppose that the coefficients $b$ and $\sigma$ satisfy Hypothesis \ref{hyp:coeff}. We also assume that $\zeta$ is a rough path fulfilling Hypothesis \ref{hyp:zeta}. The solution of the RDE (\ref{eq:relaxed-rde}) with initial condition $y$ is given by a flow map $x_{t}=U_{t\leftarrow s}^{\zeta}(y)$. In this context, the Jacobian of $U$ is defined by
    \[J_{t\leftarrow s}^{\zeta}(y)\cdot e := \frac{d}{d\varepsilon}\bigg|_{\varepsilon=0}U_{t\leftarrow s}^{\zeta}(y+\varepsilon e) .\]
    Then the following facts hold true for $J_{t\leftarrow s}$:
\begin{enumerate}[wide, labelwidth=!, labelindent=0pt, label=\textnormal{(\roman*)}]
\setlength\itemsep{.1in}    

\item Setting $\mathcal{V}_{s,t}^{e}=J_{t\leftarrow s}(y)\cdot e$, then $\mathcal{V}$ solves a RDE on $[s,T]$. Namely, for $t\in [s,T]$ we have 
    \begin{equation}\label{eq: Jacobian-RDE}
        \mathcal{V}_{s,t}^{e}=e + \int_{s}^{t}\left(\int_{U}Db(r,x_{r},a)\gamma_{r}(da) \right) \mathcal{V}_{s,r}^{e}\,dr +\int_{s}^{t}D\sigma (r,x_{r})\mathcal{V}_{s,r}^{e}\, d\zeta_{r}.
    \end{equation}


  \item The following estimate hold:
    \begin{equation} \label{ineq: Jacob}
        \|J_{\cdot \leftarrow s}(y)\|_{\infty;[s,T]}\leq C\exp{(CT\|\zeta\|_{\alpha}^{1/\al})} . \end{equation}
      \end{enumerate}
\end{proposition}

    \begin{equation} 
        \|J_{\cdot \leftarrow s}(y)\|_{\infty;[s,T]}\leq C\exp{(CT\|\zeta\|_{\alpha}^{1/\al})} . \end{equation}

\begin{lemma} \label{Lemma: X}
We work under the same assumptions as in Proposition \ref{prop: Flow and Jacobian}. 
For $\beta>0$ let $x^{\beta}$ be the solution to~(\ref{eq:relaxed-rde}), considered with $s=0$ and the relaxed control $\gamma^{\beta}$ in (\ref{spike}). Also, recall from Definition~\ref{def:spike variation} that $(\overline{x},\overline{\gamma})$ is an optimal pair. Then the function $\beta\mapsto x^{\beta}$ is right differentiable, where $x^{\beta}$ has to be understood as an element of $\mathcal{Q}_{\zeta}^{\kappa}(\R^n)$. Moreover, the right derivative $V^{\beta}:= \partial_{\beta}x^{\beta}$ satisfies the following linear rough differential equation: For $t<t_{0}+\beta$ we have $V^{\beta}_{t}=0$ and for $t\geq t_{0}+\beta$ the path $V^{\beta}$ solves 
\begin{equation} \label{eq: RDE for derivative}
V^{\beta}_{t}=g^{\beta}+\int_{t_{0}+\beta}^{t}\lp\int_{U}Db(r,x_{r}^{\beta},a)\ga_{r}^{\beta}(da)\rp V_{r}^{\beta}dr +\int_{t_{0}+\beta}^{t}D\sigma(r,x_{r}^{\beta})V_{r}^{\beta}d\zeta_{r},
\end{equation}
where the quantity $g^{\beta}$ is given by
\begin{equation} \label{eq: g-beta}
g^{\beta}=\int_{U}b(t_{0}+\beta,x^{\beta}_{t_{0}+\beta},a)\mu(da)-\int_{U}b(t_{0}+\beta,x^{\beta}_{t_{0}+\beta},a)\overline{\ga}_{t_{0}+\beta}(da). 
\end{equation}
In addition, the following estimate holds true:
\begin{equation} \label{Ineq: derivative_of_dynamics}
\mathcal{N}\left[V^{\beta};\mathcal{Q}_{\zeta}^{\ka}([0,T];\R^m)\right]\leq C\exp\lp CT\|\zeta\|_{\alpha}^{1/\alpha} \rp .
\end{equation}
In particular, one has
\begin{equation}\|x^{\beta}-\overline{x}\|_{\infty;[0,T]}=O(\beta) .
\label{eq: Pontry-1}\end{equation}
        \end{lemma}

\begin{proof}
For a generic $\beta>0$, consider the rough differential equation
\begin{equation} \label{eq: dynamics-beta-0}
    x_{t}^{\beta}=y+\int_{0}^{t}\lp \int_{U}b(r,x_{r}^{\beta},a)\ga_{r}^{\beta}(da)\rp\,dr+\int_{0}^{t}\sigma(r,x_{r}^{\beta})\,d \zeta_{r}.
\end{equation}
Then, by Proposition \ref{ExistUniq2} there exists a unique solution $x^{\beta}$ for every given initial value $y$. So for $t\geq t_{0}$, we can consider the flow map $U_{t\leftarrow t_{0}}(\overline{x}_{t_{0}})=x_{t}^{\beta}$ (Here we use the fact that $x^{\beta}\equiv \overline{x}$ on $[0,t_{0}]$). By Proposition \ref{prop: Flow and Jacobian} we obtain that $\mathcal{V}^{e}_{t_{0}t}=J_{t\leftarrow t_{0}}(\overline{x}_{t_{0}+\beta})\cdot e$ satisfies
\[\mathcal{V}_{t_{0},t}^{e}=e+\int_{t_{0}}^{t}\lp \int_{U}Db(r,x_{r}^{\beta},a)\gamma_{r}^{\beta}(da) \rp \mathcal{V}_{t_{0},r}^{e}\,dr +\int_{t_{0}}^{t}D\sigma(r,x_{r}^{\beta})\mathcal{V}_{t_{0},r}^{e}\,d\zeta_{r}.\]
We now wish to examine the differentiability of the application $\beta\mapsto x^{\beta}$. We will divide this analysis in several steps.

\smallskip
\noindent
\textit{Step 1: Reduction of the problem.}
One way to prove the differentiability of $\beta\mapsto x^{\beta}$ is to proceed by Picard iterations. Namely, we set $x_{t}^{\beta,0}=y$. Then, we iteratively define $x^{\beta,k+1}=\Gamma(x^{\beta,k})$, with $\Gamma$ given as an application
\begin{equation*}
  \Gamma\colon \mathcal{Q}_{\zeta}^{\ka}(\R^{m})  \longrightarrow \mathcal{Q}_{\zeta}^{\ka}(\R^m) ,
\quad\text{such that}\quad
  w  \longmapsto \Gamma(w) = z \, ,
\end{equation*}
 where $z$ is defined by
 \[z_{t}=y+\int_{0}^{t}\lp \int_{U}b(r,w_{r},a)\ga_{r}^{\beta}(da)\rp\,dr+\int_{0}^{t}\sigma(r,w_{r})\,d \zeta_{r}.\]
 With this notation in hand, a possible strategy to get differentiability results in \cite[Chapter 11]{friz-victoir}  is the following:
   \begin{enumerate}[wide, labelwidth=!, labelindent=0pt, label=\textnormal{(\arabic*)}]
   \setlength\itemsep{.05in}
       \item It is well known that the Picard scheme $x^{\beta,k}$ is convergent in $\mathcal{Q}_{\zeta}^{\kappa}(\R^m)$. It converges to the solution $x^{\beta}$ of (\ref{eq: dynamics-beta-0}).
       
       \item One can also prove differentiability by induction. Namely, if one assumes that $\beta \mapsto x^{\beta,k}$ is differentiable, the problem boils down to establish differentiability of
       \begin{equation}\label{c1}
 \beta\in [0,T]  \longmapsto \Gamma(x^{\beta,k})\in\mathcal{Q}_{\zeta}^{\ka}(\R^m).
\end{equation}
In particular, it will be easy to show that $x^{\beta,k}|_{[0,t_{0}+\beta]}=x^{\beta+h,k}|_{[0,t_{0}+\beta]}$.

\item\label{it:differentiability-3} 
Once we have obtained differentiability in \eqref{c1}, together with the uniform (in $k$) bounds for the derivatives, some compactness arguments allow to take limits and prove that $x^{\beta}$ in ~(\ref{eq: dynamics-beta-0}) is differentiable.
\end{enumerate}
For sake of conciseness, we will refrain from implementing the whole strategy summarized above. We will content ourselves with giving some details about differentiability in the main step~\eqref{c1}. Moreover, we shall only focus on right derivatives in $\beta$, which are sufficient for our purposes.

Summarizing our considerations so far, our differentiation problem is now reduced to the following: Assume $\beta \mapsto x^{\beta}$ is right differentiable at $\beta$. Namely, we assume that the following limit exists in $\mathcal{Q}_{\zeta}^{\kappa}$:
\begin{equation}\label{eq: Derivative of xbeta}
    X^{\beta}:=\lim_{h\downarrow 0}\frac{x^{\beta+h}-x^{\beta}}{h}.
\end{equation}
 Next, set $z^{\beta}=\Gamma(x^{\beta})$. We wish to prove that $z^{\beta}$ is also differentiable as an element of $\mathcal{Q}_{\zeta}^{\ka}$. In addition, we want to derive an equation like (\ref{eq: RDE for derivative}) and estimates such as (\ref{Ineq: derivative_of_dynamics}) for $Z^{\beta}\equiv \partial_{\beta}z^{\beta}$. The remainder of the proof will focus on this step.

\smallskip
\noindent
\textit{Step 2: Time decomposition of the dynamics.} 
In order to get an expression for $Z^{\beta}$ defined as above, we will write the equation for $z^{\beta}$ in a more enlightening way (according to the decomposition of $\gamma^{\beta}$). Specifically, we divide the interval $[0,T]$ in three  pieces:

\begin{enumerate}[wide, labelwidth=!, labelindent=0pt, label=\textnormal{(\alph*)}]
\setlength\itemsep{.05in} 

\item \label{case: t<t_{0}}
For $t<t_{0}$, the variation on $\overline{\ga}$ has not occurred yet. Specifically, going back to (\ref{spike}), we have $\gamma^{\beta}_{t}=\overline{\gamma}_{t}$ regardless of $\beta$. Hence we have 
\begin{equation}\label{eq: dynamics-equiv}
z_{t}^{\beta}={u}_{t}, 
\quad \text{for all} \quad 
t\in [0,t_{0}],
\end{equation}
where $u_{t}=\Gamma(x^{\beta}_{t})$ is independent of $\beta$.

\item
Next, for $t\in[t_{0},t_{0}+\beta]$, the dynamics (\ref{eq: dynamics-beta-0}) can be read as
\begin{equation}\label{eq:decomp-dynamics}
    z_{t}^{\beta}={u}_{t_{0}}+\int_{t_{0}}^{t}\lp\int_{U}b(r,x_{r}^{\beta},a)\mu(da) \rp\,dr +\int_{t_{0}}^{t}\sigma(r,x_{r}^{\beta})\,d\zeta_{r} .
\end{equation}
Moreover, for $h>0$ the path $y^{\beta+h}$ also follows the dynamics (\ref{eq:decomp-dynamics}) on $[t_{0},t_{0}+\beta]$.

\item
Eventually, for $t\geq t_{0}+\beta$, equation (\ref{eq: dynamics-beta-0}) becomes
\begin{multline}\label{eq:decomp-dynamics-2}
z_{t}^{\beta}={u}_{t_{0}}+\int_{t_{0}}^{t_{0}+\beta}\lp \int_{U}b(r,x_{r}^{\beta},a)\mu(da) \rp\,dr \\
+\int_{t_{0}+\beta}^{t}\lp \int_{U}b(r,x_{r}^{\beta},a)  \overline{\ga}_{r}(da) \rp \,dr 
+ \int_{t_{0}}^{t}\sigma(r,x_{r}^{\beta})\,d\zeta_{r}.
\end{multline}
\end{enumerate}
We are now ready to derive an expression for $Z^{\beta}=\partial_{\beta}z^{\beta}$, following the decompositions $(\ref{eq: dynamics-equiv})$-$(\ref{eq:decomp-dynamics})$-$(\ref{eq:decomp-dynamics-2})$ . We deal with the previous three cases separately.

\begin{enumerate}[wide, labelwidth=!, labelindent=0pt, label=\emph{(\roman*)}]
\setlength\itemsep{.1in}     

\item \emph{Case $t\in [0,t_{0})$.}
    Earlier, we saw in (\ref{eq: dynamics-equiv}) that $z^{\beta}_{t}={u}_{t}$. Similarly, we have $z^{\beta +h}_{t}={u}_{t}$. Therefore, we conclude that $Z^{\beta}_{t}\equiv 0$.
    
    \item \emph{Case $t\in [t_{0},t_{0}+\beta]$.}
Looking at (\ref{eq:decomp-dynamics}) and observing that $z^{\beta+h}_{t}$ also satisfies (\ref{eq:decomp-dynamics}), we conclude that $z^{\beta}\equiv z^{\beta+h}$. Ultimately, we yet again have $Z^{\beta}_{t}=0$.
    
    \item \emph{Case $t\in (t_{0}+\beta,T]$.}
   Invoking the dynamics of $x^{\beta}$ and $x^{\beta+h}$, and using the observation that $x^{\beta+h}_{t_{0}+\beta}=x^{\beta}_{t_{0}+\beta}$, we obtain:
   \begin{multline*}
       \frac{z_{t}^{\beta+h}-z^{\beta}_{t}}{h}
       =\frac{1}{h}\int_{t_{0}+\beta}^{t}\int_{U}b(r,x_{r}^{\beta+h},a)\ga_{r}^{\beta+h}(da)\,dr 
       -\frac{1}{h}\int_{t_{0}+\beta}^{t}\int_Ub(r,x_{r}^{\beta},a)\ga_{r}^{\beta}(da)\,dr\\
       +\frac{1}{h}\int_{t_{0}+\beta}^{t}\{\sigma(x_{r}^{\
       \beta+h})-\sigma(r,x_{r}^{\beta})\}\,d\zeta_{r}.
\end{multline*}
Therefore we get the following decomposition for the increments of $z^{\beta}$:
\begin{equation}
    \frac{z^{\beta+h}_{t}-z^{\beta}_{t}}{h}=I_{h}+II_{h}+III_{h},
\label{eq: decomp of y}
\end{equation}
where the quantity $I_{h}$ is given by
\begin{equation}\label{eq: Ih}
 I_{h} 
 =\frac{1}{h}\int_{t_{0}+\beta}^{t}\int_{U}\{b(r,x_{r}^{\beta+h},a)-b(r,x_{r}^{\beta},a)\}\ga_{r}^{\beta}(da)\,dr ,
\end{equation}
and where $II_{h}$ and $III_{h}$ are respectively defined by
\begin{align}
  \label{eq: IIh} II_{h}&=\frac{1}{h}\int_{t_{0}+\beta}^{t}\int_{U}b(r,x_{r}^{\beta+h},a)\{\ga^{\beta+h}_{r}-\ga^{\beta}_{r}\}(da)\, dr\\
     \label{eq: IIIh}  III_{h}&=\frac{1}{h} \int_{t_{0}+\beta}^{t}\{\sigma(r,x_{r}^{\beta+h})-\sigma(r,x_{r}^{\beta})\}\,d\zeta_{r}
\end{align}
We shall derive the limits in those three terms separately.


\smallskip
\noindent
\textit{Step 3: Limit for $I_{h}$.} 
Let us deal with the term $I_{h}$ in (\ref{eq: Ih}) first. This is the easiest of the terms in~\eqref{eq: decomp of y}. Indeed, since we have assumed $x^{\beta}$ is differentiable in~(\ref{eq: Derivative of xbeta}), one can easily check that
\begin{equation}
\lim_{h\downarrow0}I_{h}=\int_{t_{0}+\beta}^{t}\lp\int_{U}Db(r,x_{r}^{\beta},a)\ga_{r}^{\beta}(da) \rp X_{r}^{\beta}\,dr.
\label{eq: limit Ih}\end{equation}


\smallskip
\noindent
\textit{Step 4: Limit for $III_{h}$.} 
We now handle the term $III_{h}$, that we wish to differentiate inside the rough integral in~\eqref{eq: IIIh}. We claim that 
\begin{equation}\label{eq: Integral_3}
    \lim_{h\to 0}III_{h}=\int_{t_{0}+\beta}^{t} D\sigma(r,x_{r}^{\beta})X_{r}^{\beta}\,d\zeta_{r} .
\end{equation}
In order to prove (\ref{eq: Integral_3}), let us give a more explicit expression for the rough integral. Namely (see \cite{gubinelli} for details) for $x^{\beta}$, the solution to (\ref{eq: dynamics-beta-0}), the strongly controlled decomposition for the rough integral can be written
\begin{equation}\label{eq: Rough_Exp-Int}
    \int_{u}^{v}\sigma(r,x_{r}^{\beta})\, d\zeta_{r}=\sigma(u,x_{u}^{\beta})\delta \zeta_{uv}+ D\sigma(u,x_{u}^{\beta})\sigma(u,x_{u}^{\beta})\zeta_{uv}^{2}+ \Lambda(\ell^{\beta})_{uv},
\end{equation}
where $\Lambda$ is the sewing map introduced in Proposition \ref{prop:La} and $\ell^{\beta}$ is an increment in $\mathcal{C}_{3}^{3\ka}$ defined by
\begin{equation}\label{eq: ell-beta}
    \ell^{\beta}=\delta\lp \sigma(\cdot,x^{\beta})\delta \zeta +D\sigma(\cdot,x^{\beta})\sigma(\cdot,x^{\beta})\zeta^{2}\rp.
\end{equation}
Now if $x^{\beta}$ is differentiable in $\beta$, it is clear that the first two terms in the right hand side of~(\ref{eq: Rough_Exp-Int}) are also differentiable. The term $\ell^{\beta}$ as given by (\ref{eq: ell-beta}) is also clearly differentiable. In order to see that $\Lambda(\ell^{\beta})$ is differentiable, we use 
(\ref{eq: ell-beta}) in Proposition \ref{prop:La} to assert
\begin{align*}
   \left\| \Lambda(\ell^{\beta+h})- \Lambda(\ell^{\beta})-h\Lambda(\partial_{\beta}\ell^{\beta})\right\|_{3\kappa}
   &=
   \left\| \Lambda(\ell^{\beta + h} -\ell^{\beta}-h\partial_{\beta}\ell^{\beta}) \right\|_{3\kappa} \\ 
   &\lesssim 
   \left\|\ell^{\beta+h}-\ell^{\beta}-h\partial_{\beta}\ell^{\beta}\right\|_{3\kappa}=o(h) \, ,
\end{align*}
where the last identity stems from the fact that $\ell^{\beta}$ is differentiable with respect to $\beta$. This argument implies that $\partial_{\beta}\Lambda(\ell^{\beta})=\Lambda(\partial_{\beta}\ell^{\beta})$. Plugging this identity back into (\ref{eq: Rough_Exp-Int}), we have proved that (\ref{eq: Integral_3}) holds true.

\smallskip
\noindent
\textit{Step 5: Limit for $II_{h}$.} 
Let us now turn to the analysis of the term $II_{h}$ in (\ref{eq: IIh}). Owing to the expressions $\ga^{\beta}$ and $\ga^{\beta+h}$, it is readily checked that 
\[II_{h}=\frac{1}{h}\int_{t_{0}+\beta}^{(t_{0}+\beta+h)\wedge t}\int_{U}b(r,x_{r}^{\beta+h},a)\{\mu-\overline{\ga}_{r}\}(da)\,dr.\]
We will further decompose this expression as 
\begin{equation}\label{eq:IIh limit}
    II_{h}=II_{h}^{(1)}+II_{h}^{(2)},
\end{equation}
where $II_{h}^{(1)}$, $II_{h}^{(2)}$ are respectively defined by
\begin{align}
    II_{h}^{(1)}&=\frac{1}{h}\int_{t_{0}+\beta}^{(t_{0}+\beta+h)\wedge t}\int_{U}\{b(r,x_{r}^{\beta+h},a)-b(r,x_{r}^{\beta},a)\}\{\mu-\overline{\ga}_{r}\}(da)\,dr\\
    II_{h}^{(2)}&=\frac{1}{h}\int_{t_{0}+\beta}^{(t_{0}+\beta+h)\wedge t}\int_{U}b(r,x_{r}^{\beta},a)\{\mu-\overline{\ga}_{r}\}(da)\,dr.
\end{align}
Moreover, if $\beta\mapsto x^{\beta}$ is differentiable and $b$ is smooth, we have
\[|b(r,x_{r}^{\beta+h},a)-b(r,x_{r}^{\beta},a)|\lesssim|h|,\]
uniformly in $(r,a)$. Therefore, combining this regularity with the regularity given by the time integral in $II_{h}^{(1)}$, we trivially  get
\begin{equation}\label{eq: II-1h limit}
    \lim_{h\downarrow 0}II_{h}^{(1)}=0.
\end{equation}
As far as $II_{h}^{(2)}$ is concerned, one can apply the fundamental theorem of calculus and obtain
\begin{equation}\label{eq: II-2h limit}
    \lim_{h\downarrow 0}II_{h}^{(2)}= \int_{U}b(t_{0}+\beta,x^{\beta}_{t_{0}+\beta},a)\mu(da)-\int_{U}b(t_{0}+\beta,x^{\beta}_{t_{0}+\beta},a)\overline{\ga}_{t_{0}+\beta}(da).
\end{equation}
Gathering (\ref{eq: II-1h limit}) and (\ref{eq: II-2h limit}) into (\ref{eq:IIh limit}), we end up with
\begin{equation}\label{eq: limit IIh}
    \lim_{h\downarrow 0} II_{h}=\int_{U}b(t_{0}+\beta,x^{\beta}_{t_{0}+\beta},a)\mu(da)-\int_{U}b(t_{0}+\beta,x^{\beta}_{t_{0}+\beta},a)\overline{\ga}_{t_{0}+\beta}(da).
\end{equation}

\smallskip
\noindent
\textit{Step 6: Proof of differentiability.} \label{lem 3.8: step 6}
We can now summarize and obtain an expression for $Z^{\beta}=\partial_{\beta}z^{\beta}$. Namely putting together (\ref{eq: limit Ih}), (\ref{eq: Integral_3}) and (\ref{eq: limit IIh}), and recalling our decomposition (\ref{eq: decomp of y}), we obtain that $Z^{\beta}$ satisfies

\begin{multline}\label{eq: Equation-Derivative 1}
Z^{\beta}_{t}=\int_{t_{0}+\beta}^{t}\lp\int_{U}Db(r,x_{r}^{\beta},a)\ga_{r}^{\beta}(da) \rp X_{r}^{\beta}\,dr + \int_{t_{0}+\beta}^{t}D\sigma(r,x_{r}^{\beta})X_{r}^{\beta}\,d\zeta_{r}\\
    \hspace{0.25cm}+ \int_{U}b(t_{0}+\beta,x^{\beta}_{t_{0}+\beta},a)\mu(da)-\int_{U}b(t_{0}+\beta,x^{\beta}_{t_{0}+\beta},a)\overline{\ga}_{t_{0}+\beta}(da).
\end{multline}
Here again, it should be mentioned that the limits in $h$ have to be understood in the rough paths sense.
\end{enumerate}

Recall that in point \ref{it:differentiability-3} of our strategy for differentiation, we took limits in Picard iterations related to~(\ref{eq: Equation-Derivative 1}). Skipping the tedious details for the justification of this step, we obtain that if $x^{\beta}$ is the solution of (\ref{eq: dynamics-beta-0}), and if we set $V^{\beta}=\partial_{\beta}x^{\beta}$, then $V^{\beta}$ solves the following linear rough differential equation on $[t_{0}+\beta,T]$:
\begin{multline}\label{eq: Equation-Derivative 2}
V^{\beta}_{t}=\int_{t_{0}+\beta}^{t}\lp\int_{U}Db(r,x_{r}^{\beta},a)\ga_{r}^{\beta}(da) \rp V_{r}^{\beta}\,dr + \int_{t_{0}+\beta}^{t}D\sigma(r,x_{r}^{\beta})V_{r}^{\beta}\,d\zeta_{r}\\
    \hspace{0.25cm}+ \int_{U}b(t_{0}+\beta,x^{\beta}_{t_{0}+\beta},a)\mu(da)-\int_{U}b(t_{0}+\beta,x^{\beta}_{t_{0}+\beta},a)\overline{\ga}_{t_{0}+\beta}(da).
\end{multline}
We have thus proved relation (\ref{eq: RDE for derivative}).

\smallskip
\noindent
\textit{Step 7: Bounds on the derivative.} 
We now turn to the proof of inequality (\ref{Ineq: derivative_of_dynamics}). Here note that we are looking for a bound for the whole norm $\mathcal{N}\left[ V^{\beta};\mathcal{Q}_{\zeta}^{\kappa} \right]$, as defined in (\ref{eq:norm-controlled}). For the sake of conciseness, we will just bound $\mathcal{N}[V^{\beta}; \mathcal{C}_{1}^{\infty}(\R^m)]$, leaving the other lengthy computations to the reader. Now in order to estimate the sup norm of $V^{\beta}$, we resort to the fact that $V^{\beta}\equiv 0$ on $[0,t_{0}+\beta)$. Furthermore, on $[t_{0}+\beta,T]$, the path $V^{\beta}$ satisfies the same equation as $J^{\zeta}_{t\leftarrow t_{0}+\beta}(y)\cdot \hat{e}$, where $J$ stands for the Jacobian in Proposition \ref{prop: Flow and Jacobian} and $\hat{e}$ is given by
\begin{equation}
   \hat{e}= \int_{U}b(t_{0}+\beta,x^{\beta}_{t_{0}+\beta},a)\mu(da)-\int_{U}b(t_{0}+\beta,x^{\beta}_{t_{0}+\beta},a)\ga_{t_{0}+\beta}(da) .
\end{equation}
Therefore, invoking the estimate \eqref{ineq: Jacob} on the Jacobian and the fact that $b$ (and thus $\hat{e}$) is uniformly bounded, we obtain
\begin{multline*}
   \mathcal{N}[V^{\beta};\mathcal{C}_{1}^{0}([0,T])]=\sup_{t\in[t_{0}+\beta,T]}|V^{\beta}_{t}|
    =\sup_{t\in[t_{0}+\beta,T]}|J_{t\leftarrow t_{0}+\beta}^{\zeta}(y)\cdot \hat{e}|\\
   \leq C_{b}\sup_{t\in[t_{0}+\beta,t]}\|J_{t\leftarrow t_{0}+\beta}(x^{\beta}_{t_{0}+\beta})\|_{\mathrm{op}}
    \leq C\exp\lp CT\|\zeta\|_{\alpha}^{1/\alpha} \rp,
    \end{multline*} 
    which proves
    \begin{equation}\|V^{\beta}\|_{\infty;[0,T]}\leq C\exp \lp C T \|\zeta\|_{\alpha}^{1/\alpha} \rp.
\label{eq: uniform-estimate}    \end{equation}
Note again that we have focused on the uniform norm in (\ref{eq: uniform-estimate}), letting the reader complete the other estimates leading to (\ref{Ineq: derivative_of_dynamics}).

Eventually, observing that $x^{0}\equiv \overline{x}$, we get that 
\begin{equation}\label{eq: FTC x}x^{\beta}_{t}-\overline{x}_{t}=\int_{0}^{\beta}V_{t}^{u} \, du.\end{equation}
Hence, relation (\ref{eq: uniform-estimate}) directly yields (\ref{eq: Pontry-1}). This finishes our proof.\end{proof}

\begin{thebibliography}{99}
\bibitem[CHT]{CHT} P.~Chakraborty, H.~Honnappa, and S.~Tindel. \newblock Pathwise relaxed optimal control of rough differential equations. \newblock \emph{arXiv preprint arXiv:2402.17900}, 2024.
\bibitem[friz-v]{friz-victoir} P.~K. Friz and N.~B. Victoir. \newblock \emph{Multidimensional stochastic processes as rough paths: theory and applications}, volume 120. \newblock Cambridge University Press, 2010.
\bibitem[gubine]{gubinelli} M.~Gubinelli. \newblock Controlling rough paths. \newblock \emph{Journal of Functional Analysis}, 216\penalty0 (1):\penalty0 86--140, 2004.
\end{thebibliography}
\end{document}
